% ═══════════════════════════════════════════════════════════════════
%  Conclusion and Future Work
% ═══════════════════════════════════════════════════════════════════

\label{sec:conclusion}

\subsection{Summary}

This thesis set out to investigate whether LLM agents produce economically interpretable negotiation behaviour in a structured bilateral market, and to examine how simulator design choices---feasibility enforcement, prompt architecture, information structure---shape the outcomes that experiments can observe and attribute to agents.

The central empirical finding is a micro--macro capability gap: LLM agents are effective deal-closers but poor price-discovery mechanisms.  Market liquidity reaches 81\% without gate assistance, individual sessions settle reliably within the round limit, and deal success rates range from 97.2\% in fixed-scenario experiments to 66.7\% under adverse market conditions.  Yet price elasticities to fundamental shocks are only 0.035--0.275, and the Experiment~4 extension shows that mean transaction prices span approximately \$30 across anchor modes under identical fundamentals.  Agents anchor to the reference price signal in the prompt rather than to their private constraints, producing strong deal convergence at informationally inaccurate prices.

Supporting this finding, Experiments~1 and~2 document concave (Boulware-style) concession trajectories with a 5.65:1 buyer-to-seller concession ratio and deadline-driven settlement acceleration (82.4\% of deals in the final two rounds; 25\% higher concession rate under a six-round horizon).  These micro-level patterns are directionally consistent with human bargaining literature.  At the market level, prices exhibit persistent within-tick dispersion without convergence, directional but attenuated shock responses, and a demand-supply asymmetry that runs counter to classical supply-demand theory.  Surplus-maximising matching improves welfare by 7.2\% relative to random matching, demonstrating that mechanism design can improve aggregate outcomes independently of agent sophistication.

A consistent methodological insight runs through the programme: protocol-level design choices---including the reference price framing, the accept-condition block, and the finite-horizon counteroffer action space---shape the behaviours experiments can observe.  All findings should therefore be understood as characterising a specific model-plus-prompt-plus-protocol combination, not LLM negotiation capability in general.

\subsection{Implications}

The findings carry implications for four audiences.

\paragraph{For understanding AI negotiation behaviour.}
Under the tested conditions, LLM agents reproduce the directional qualitative signatures of deadline pressure, Boulware concession, and first-offer anchoring documented in human negotiation studies.  These alignments suggest viability for testing protocol designs and studying how informational environments shape bargaining outcomes.  However, systematic departures---price stickiness and the absence of threat-of-rejection dynamics under the counteroffer protocol---limit the extent to which LLM agents mimic the full range of human negotiation behaviour.  Direct human comparison under identical protocol is required before stronger claims can be made.

\paragraph{For experimental methodology.}
Simulator design choices---feasibility enforcement, prompt steering, reference price framing---can mask or distort the behaviours experiments are designed to study.  The gate-off design was motivated precisely by preliminary observations that gate-on runs suppressed deadline effects and concession trajectories; the Experiment~4 extension provides a parallel lesson for reference price framing.  Researchers using LLM agents in structured settings should treat such interventions as controllable variables, report their involvement rates, and conduct ablation studies to separate agent capability from system-level enforcement.

\paragraph{For automated negotiation deployment.}
LLM agents reach agreement in the majority of sessions when the ZOPA is wide (overall deal rate: 78.6\%), but price stickiness limits their usefulness as price-discovery mechanisms.  Deployed systems using similar prompt architectures may exhibit systematic mis-pricing that does not self-correct as market conditions change, because agents anchor to the reference price embedded in the prompt rather than adapting to new fundamentals.  The scope of this concern beyond the tested model--prompt--protocol combination is unknown.

\paragraph{For economic theory.}
Several patterns---price dispersion, first-mover advantage at the micro level, and directional shock response---are directionally consistent with bilateral bargaining theory, even though agents lack sophisticated strategic reasoning.  The current experiments cannot determine whether these patterns arise from the bargaining structure, model-specific training distributions, or prompt architecture; no experiment independently varies agent sophistication.  The appropriate interpretation is modest: qualitative predictions of bilateral bargaining theory can emerge in a simulator with unsophisticated agents, not that LLM agents are strategically rational or constitute valid substitutes for human participants in experimental economics.

\subsection{Future Work}

Several directions would strengthen and extend the current findings.

\paragraph{Prompt ablation.}
Prompt-level ablation---removing the accept-condition block, removing the monotonic concession instruction, extending the action space with a voluntary walk-away option, and varying the deadline-salience language---is needed to clarify the boundary between agent-driven behaviour and protocol-induced compliance.  the Experiment~4 extension demonstrates that reference price framing is a substantial source of price-level variation; similar ablations for other prompt components would clarify how much of the observed bargaining behaviour reflects genuine LLM decision-making.  The platform already supports prompt configuration via YAML, making these ablations straightforward to execute.

\paragraph{Multi-model comparison.}
All reported experiments use a single model (\texttt{Qwen2.5-14B-Instruct}).  Repeating the experimental configuration with models of different sizes (7B, 70B) and architectures under identical prompt and gate settings would test whether price stickiness, asymmetric concession burden, and deadline sensitivity are model-scale phenomena or general properties of LLM bargaining.  This is the highest remaining external-validity gap.

\paragraph{Gate ablation.}
A within-configuration gate ablation---comparing gate-on and gate-off under the same model---would provide controlled evidence on how feasibility enforcement affects observed outcomes.  This would complement the preliminary observation that gate-on runs suppressed behavioural signals.

\paragraph{Extended market replication.}
Experiments~3 and~4 were replicated across three seeds, which improves confidence in the directional findings but reveals substantial cross-seed variance (notably in Experiment~4's shock response trajectories).  Additional seeds and longer tick horizons (beyond 10 ticks) would help distinguish persistent market-level patterns from seed-specific fluctuations.

\paragraph{Supply-demand mechanism isolation.}
Experiment~6 establishes the asymmetric price response to symmetric cost and value shifts but cannot isolate the responsible mechanism.  Three targeted ablations would address this: (i)~varying the size of the demand shift to test whether the price response is linear or threshold-gated; (ii)~running the supply-shock condition with different first-mover protocols (seller-first or simultaneous) to test whether cost-floor pass-through depends on the buyer-first turn order; (iii)~testing with and without the reference price signal to isolate anchoring effects.

\paragraph{Reputation and communication experiments.}
The platform supports two additional experiment families---reputation and repeated interaction (using the round-robin matcher and reputation agent) and communication strategy ablation (using prompt-level communication modifiers)---for which infrastructure is implemented but empirical evaluation has not been conducted.  Both require the gate to be disabled so that strategic adaptation can manifest over multiple rounds.

\paragraph{Statistical analysis.}
The post-hoc anchor-mode comparison (Section~\ref{sec:results-anchor}) represents a first step toward formal statistical evaluation, but most findings in this thesis remain at the descriptive level.  Broader application of hypothesis tests, confidence intervals, and effect-size measures---together with additional seeds to increase independent replication---would strengthen the evidential basis of the current findings.

\paragraph{Human--LLM comparison.}
Running the same experimental conditions with human participants would enable direct comparison of bargaining behaviour.  This would clarify whether the observed patterns are specific to LLMs or whether they extend to any agent operating with explicit constraint information.

\paragraph{Stronger baselines.}
The rule-based agent provides a deterministic baseline but follows a simple linear concession strategy.  Adding stronger algorithmic baselines---such as Nash-equilibrium-computing agents or game-theoretic optimal strategies---would help calibrate the gap between LLM behaviour and theoretically optimal play.

\vspace{1em}

The platform is a testbed, not a finished product.  Its value lies in enabling structured, logged comparison of LLM bargaining behaviour against theoretical reference points and documented human negotiation patterns---and in making simulator design choices, such as the informational anchor environment and prompt architecture, explicit and configurable variables rather than invisible background conditions.
